\section{Introduction}

Recent research on small and medium-sized enterprise (SME) resilience emphasizes that digitalization, innovation, and organizational learning increasingly shape how firms absorb disruption and adapt under crisis conditions \citep{sagala2024antifragility, sinha2025digitalisation}. This contemporary concern builds on foundational development-economics evidence that SMEs are important engines of employment, innovation diffusion, and inclusive growth in developing economies \citep{beck2006small, ayyagari2007small}. In Rwanda, as in many emerging economies, SMEs operate within a changing policy and market environment shaped by digital transformation, financial constraints, and recurrent demand shocks. These pressures make early identification of firm vulnerability a central policy concern, but realized closure or exit is often observed too late for preventive support.

This paper reframes SME vulnerability as a predictive fragility-screening problem. Rather than estimating the causal effect of a single intervention or predicting observed survival duration, we ask whether a constructed fragility classification can be reproduced, audited, and predicted transparently from enterprise-survey variables. Four questions organize the analysis: Can the legacy SME fragility classification be reconstructed from raw survey variables? To what extent does full-model performance reflect outcome-rule replication? How much independent predictive signal remains after removing outcome-rule variables? What governance conditions are required before such a score can inform policy? The distinction matters. A predictive model can support screening and prioritization even when cross-sectional data cannot identify causal effects \citep{shmueli2010explain}.

The empirical setting is the Rwanda 2023 Enterprise Survey. The project folder contained raw survey data, curated predictor files, Stata-derived classification exports, draft manuscripts, and model figures. We reorganized these materials into a reproducible workflow and generated new Python and Stata outputs from the raw and curated analysis files. The resulting analysis uses 358 formal firms covered by the World Bank Enterprise Survey frame. The data should not be interpreted as covering all Rwandan SMEs, informal enterprises, microenterprises, or firms below the Enterprise Survey eligibility threshold.

The paper makes three contributions. First, it presents a transparent firm-level fragility-screening framework suitable for social-science research using cross-sectional enterprise data in a low-data African enterprise context. Second, it demonstrates why predictive performance must be interpreted alongside outcome construction: the full model validates the mechanics of the recovered index, while the leakage-reduced model provides the more substantively meaningful benchmark for independent predictive signal. Third, it provides a generalizable workflow for auditing constructed policy-screening outcomes in settings where index labels, variable coding, and predictive performance could otherwise produce misleading policy claims. The workflow can inform research on enterprise surveys, social registries, credit-screening indices, and development-program targeting tools when the outcome is constructed rather than directly observed.
